% Section 17.10, from lectures/L17/appendix/b_exercises.md.

\section{Exercises}
\label{c17:sec:exercises}

\begin{aside}
Once your receive logic is finished, \file{controller/can_controller_tb.vhd} verifies transmission and
reception together, and \code{make build-project} stops skipping it: all six testbenches in
\file{controller/} run from here on. The two in \file{bridge/} go on being skipped until
\chapref{c18:ch} and \ref{c19:ch} write the modules they exercise.

If it is still reported as skipped after you believe the receive path is finished, read the message
before your design. The build settles the matter by looking for \chapref{c17:ch}'s receive-side CRC
gating inside \file{can_controller.vhd}, so a correct controller spelling the two expressions
differently is skipped rather than run. \code{CI_BUILD_ALL=1 make build-project} forces it, and
appendix~\ref{app:sim} explains why the check works that way.
\end{aside}

\exercise{Reasoning about the RX-side timing fix}{Reasoning}
\label{c17:ex:timing}
\xpart{a} \Secref{c17:sec:preempt} shows that \code{rx_shift_reg} has to be enabled
\emph{combinationally}, during the \code{STATE_LOAD_*} state itself, instead of waiting until
\code{state} visibly becomes the shift state. Explain in your own words why first waiting for the
state transition to complete would make \code{rx_shift_reg} miss the very sample that triggered the
wait in the first place.

\xpart{b} \code{STATE_CRC_LO_WAIT} waits for \code{bt_bit_done}, not for \code{bt_sample}, unlike
every wait at a group boundary in the design (the \code{STATE_LOAD_*} states). Explain in your own
words why that transition in particular has to align with a \emph{bit period boundary} and not with a
sample point. Which other states share the \code{bt_bit_done} trigger, and what do they have in
common with it?

\xpart{c} Suppose (hypothetically) that \code{STATE_CRC_WAIT} also waited for \code{bt_sample} in the
RX role, on top of \code{STATE_LOAD_CRC_HI} doing the same immediately afterwards. Explain what would
go wrong, and why it would show up not so much as an outright failure as a receiving node ending up
one bit period further behind a transmitting node than it should be.

\exercise{Arbitration by hand}{Reasoning}
\label{c17:ex:arbhand}
Two nodes request transmission on the same cycle: node A with ID \code{0x2A0}, node B with ID
\code{0x2C0}.

\xpart{a} Write both IDs in binary (11 bits, MSB first) and mark the first bit position where they
differ.

\xpart{b} Using \secref{c10:sec:arbitration}'s arbitration rule (dominant \code{0} wins), say which
node wins and at which bit the loser drops out.

\xpart{c} The losing node detects its loss with
\code{bt_sample = '1' and txsr_tx_bit = '1' and rx_bus_s2 = '0'}. For the loser at that bit, say what
each of the three terms evaluates to and why the condition fires.

\xpart{d} The loser sets \code{error <= '1'} and drops back to \code{STATE_IDLE} while the winner is
still driving the bus. \code{error} is cleared in \code{STATE_START}. Explain why the loser's
\code{error} would become unobservable to a caller if \code{STATE_IDLE} entered a reception on a
dominant \emph{level}.

\xpart{e} A tempting fix is to look for the \emph{edge} from recessive to dominant instead, in line
with how \secref{c10:sec:frameformat} describes SOF. Show that that does not really work either,
using the bit stuffing rule: what does the transmitter insert after five dominant bits in a row, and
what does the bus do on the bit after that? Roughly how far into the winner's frame does the loser get
before the edge test fires anyway?

\xpart{f} State the rule the design actually uses (\secref{c16:sec:roles}). Two facts make it hold up;
give both: what the bit stuffing rule bounds inside the stuffed region, and what the unstuffed tail
adds on top of that which pushes the threshold to eleven. Then say why the idle counter has to be
cleared to \textbf{full} and not to empty when the node comes out of reset.

\exercise{The SOF bit nobody keeps}{Reasoning}
\label{c17:ex:sof}
\code{STATE_SOF} enables \code{rx_shift_reg} for one bit and then throws the value away.
\Secref{c17:sec:states} explains why; this exercise is about seeing the failure with your own eyes.

\xpart{a} Take ID \code{0x000}, DLC \code{0}. Write out the bits a transmitter puts on the wire from
SOF through the end of the control field, using \secref{c16:sec:fields}'s field table, and mark where
\code{tx_shift_reg} inserts every stuff bit. Remember that the run counter counts SOF.

\xpart{b} Now do the same as seen by a receiver that skipped SOF, so that its run counter is one
dominant bit behind. Mark where \emph{it} expects stuff bits. Line the two up against each other and
state the first bit position where they disagree.

\xpart{c} At that position the transmitter sends a stuff bit and the receiver expects none. What does
the receiver do with it, and what does that make of every bit after it?

\xpart{d} Repeat \textbf{a)} and \textbf{b)} for ID \code{0x123}, the ID \code{can_controller_tb} uses
in case 1, and confirm that the two counters agree all the way. State the general rule for which
identifiers break: what has to be true of the first identifier bits, and how many of the 2048 standard
IDs satisfy it?

\xpart{e} Case 2 in \code{can_controller_tb} \emph{does} send ID \code{0x000}. Explain why the bug
would still not show up there, and say what you would add to the testbench to catch it.

\exercise{Receiving a frame with no data}{Reasoning}
\label{c17:ex:dlc0}
Walk through the \textbf{receive role} over a complete frame with \code{DLC = 0}, using the per-state
table in \secref{c17:sec:states}. A receiver does not know the DLC until it has decoded the control
field, so this is also the path that tests whether your transition in \code{STATE_CTRL} is right.

\xpart{a} List every state the node passes through, from \code{STATE_IDLE} back to
\code{STATE_IDLE}. Mark the one transition that differs from a frame carrying data, and say what makes
it fire.

\xpart{b} For every \code{STATE_LOAD_*} state in your list, state the \code{rxsr_bit_count} the
combinational network pre-empts, and say which \code{bt_sample} that pre-emption has to be in place
for: the one triggering the transition out of \code{STATE_LOAD_*}, or the next one.

\xpart{c} Suppose the pre-emption were one cycle too late everywhere, that is, that \code{rxsr_enable}
went high only once \code{state} visibly read the shift state. How many real bits would
\code{rx_shift_reg} have missed for a frame with \code{DLC = 0} by the time the node reaches
\code{STATE_CRC_LO_WAIT}, and what would \code{crc_valid} read there?

\xpart{d} Your answer to \textbf{c)} is a CRC error reported many bit periods after the actual
mistake. Relate that to \secref{c17:sec:preempt}'s claim that a one-cycle misalignment accumulates
into a silent misalignment that only shows up later, far from its cause, and say why that makes this
kind of bug expensive to debug from the testbench's output alone.
